\itemhead{1.2 例题 · 例题8（教材p.21）}{XJ-B1-C01-EX-P021-Y0177}
\label{XJ-B1-C01-EX-P021-Y0177}
{\footnotesize 湘教版教材p.21 / PDF第28页 / 1.2 例题 / 例题8}
\field{原题}{判断下列存在量词命题的真假：
\subq{1}$\exists a\in\mathbb Z,\ a^2=3a-2$；
\subq{2}$\exists a\ge3,\ a^2=3a-2$；
\subq{3}设 $A,B,C$ 是平面上不在同一直线上的三点，在该平面上存在某个点 $P$，使得 $PA=PB=PC$。}
\field{参考答案}{（1）真，$a=1$ 或2；（2）假，方程仅有根1、2；（3）真，三角形三边垂直平分线的交点即外心。}
{\small \field{关联知识}{L03 全称、存在量词及真假的判定（教材帮p.38、39 / PDF47、48）}
\field{方法与易错点}{先确认量词、作用对象及论域，再采用任意性证明或实例反例。 易错：用几个满足的数证明全称命题；把“存在”理解为“唯一存在”。}
\field{答案依据}{例题原解位置：湘教版教材p.21及题库原图；独立编写可编辑参考解答。答案册此部分主要收练习、习题，不伪造例题答案册位置。}}
